\clearpage\thispagestyle{empty}
\pdfbookmark[0]{Source for the 27 June 1808 Letter to von Lindenau}{lindenau1808june-sources}
\begin{center}\textbf{Source for the 27 June 1808 Letter to von Lindenau}\end{center}
\noindent This AI translation follows Carl Friedrich Gauss’s letter to Bernhard August von Lindenau of 27 June 1808 in \href{https://gdz.sub.uni-goettingen.de/id/PPN236007467}{Gauss’s \textit{Werke}, vol.~VI} (Göttingen, 1874), pp.~301--303. It first appeared as item~IX, \textit{Auszug aus einem Schreiben des Hrn. Professors Gauß}, in the July 1808 issue of \textit{Monatliche Correspondenz}, vol.~XVIII, \href{https://www.digitale-sammlungen.de/en/view/bsb10538609?page=95}{pp.~83}, \href{https://www.digitale-sammlungen.de/en/view/bsb10538609?page=96}{84}, \href{https://www.digitale-sammlungen.de/en/view/bsb10538609?page=97}{85} and \href{https://www.digitale-sammlungen.de/en/view/bsb10538609?page=98}{86}, in the digitized copy of the Bavarian State Library.
\par\medskip
The first printing gives Harding’s Dollond instrument as $3\tfrac12$ feet. The collected edition gives $3\tfrac14$ feet. The tables retain the collected readings, including Vesta’s aphelion in the letter and the perihelion in the separate report of 4 July.
\par\medskip
Gauss describes the four-observation problem in \textit{Theory of the Motion of the Heavenly Bodies Moving About the Sun in Conic Sections}, Book~II, Section~II. See \href{https://www2.karlin.mff.cuni.cz/~tuma/Aplikace21/theoryofmotionof00gausuoft.pdf#page=258}{Art.~164} and the \href{https://www2.karlin.mff.cuni.cz/~tuma/Aplikace21/theoryofmotionof00gausuoft.pdf#page=267}{Vesta example in Art.~171}, pp.~243--248, in Charles Henry Davis’s 1857 English translation. The book gives its own calculation and tabulated elements. Translation and typesetting produced by AI.
\clearpage
